% STATUS: DRAFT
\chapter{Almost-commutative geometries and the Standard Model}\label{ch:almost_comm}
\skippable

\begin{physmotiv}
Take spacetime times a tiny \emph{finite} noncommutative space. The Dirac operator acquires components along the finite part; gauge fields come from the manifold directions and the Higgs field from the discrete directions, as parts of a single ``connection''. With the right finite algebra, the gauge group and fermion representations of the Standard Model come out. This is the Chamseddine--Connes--Marcolli picture. It is elegant, makes predictions, and has an honest, mixed track record, which we describe.
\end{physmotiv}

\begin{unsettled}
The Standard Model reconstruction is a research program with contested status. Statements about predictions and their fate below are summarized from memory of the literature and must be verified (see \texttt{refs.bib}).
\end{unsettled}

\section{Product geometries}

An \emph{almost-commutative} geometry is the product of a manifold triple $(C^\infty(M),L^2(M,S),D_M)$ and a finite triple $(\A_F,\HH_F,D_F)$: algebra $C^\infty(M)\otimes\A_F$, Hilbert space $L^2(M,S)\otimes\HH_F$, $D=D_M\otimes\id+\gamma_5\otimes D_F$. The finite algebra $\A_F$ is a sum of matrix algebras. \emph{Inner fluctuations} $D\mapsto D+A+JAJ^{-1}$ with $A=\sum a_i[D,b_i]$ produce: gauge bosons from the $M$-components of $A$, and scalar fields (Higgs) from the $F$-components, of the unitary group of the algebra; $D_F$ encodes Yukawa couplings and masses.

\section{The finite algebra of the Standard Model}

The choice $\A_F=\C\oplus\mathbb H\oplus M_3(\C)$ (with $\mathbb H$ the quaternions) acting on a space whose basis are the fundamental fermions and their antiparticles reproduces the $SU(3)\times SU(2)\times U(1)$ structure and the hypercharges after imposing the first-order condition and a KO-dimension $6\bmod8$ finite geometry (Chamseddine--Connes). The spectral action then yields the bosonic Lagrangian: Yang--Mills terms, Higgs potential, and gravity, with relations among the couplings at the cutoff scale.

\section{What is predicted, and what is not}

\begin{itemize}
\item \textbf{Relations at the unification scale:} gauge coupling unification relations (like $g_2^2=g_3^2=\tfrac53g_1^2$) and a relation between the Higgs quartic, top Yukawa and gauge couplings. Running these down to low energies predicted a Higgs mass of order 170 GeV (Chamseddine--Connes--Marcolli). The measured 125 GeV is in tension with this prediction; proposed repairs include an additional real scalar field coupled to the Higgs and neutrino sector (Chamseddine--Connes--van Suijlekom and others) and other modifications (e.g.\ departures from the original assumptions about the Dirac operator). \textbf{Status: contested.}
\item \textbf{Number of generations and gravity:} the finite space accommodates three generations as input, not a derivation.
\item \textbf{Neutrinos:} Majorana-type terms and the see-saw can be incorporated.
\end{itemize}

\begin{whatwrong}
This is not a proof that the Standard Model ``is'' noncommutative geometry; a different finite geometry, or the failure of the unification boundary conditions, changes the predictions. The program is best seen as a constraint-satisfaction exercise: a small set of axioms (spectral triple axioms) selects a restricted family of field content and relations.
\end{whatwrong}

\section*{Exercises}
\begin{exercise}
For $\A_F=\C\oplus\C$, $\HH_F=\C^2$, $D_F=\begin{pmatrix}0&m\\m&0\end{pmatrix}$ compute the inner fluctuations $A=\sum a[D_F,b]$ and show they form a single complex scalar field $\phi$ with $D_F\to\begin{pmatrix}0&m+\phi\\\bar m+\bar\phi&0\end{pmatrix}$; compute the spectral action $\Tr f(D_F/\Lambda)$ and identify a quartic potential for $\phi$ (a toy Higgs).
\end{exercise}
\begin{exercise}
Explain how the Lie algebra of the unitary group of $\C\oplus\mathbb H\oplus M_3(\C)$ contains $\mathfrak u(1)\oplus\mathfrak{su}(2)\oplus\mathfrak u(3)$ and how unimodularity cuts $\mathfrak u(3)\to\mathfrak{su}(3)$ and constrains hypercharges.
\end{exercise}
